% =====================================================================
% Primera parte – 1. Problema y contexto
% =====================================================================
\section{Problema y contexto}

\subsection{Planteamiento del problema}

\paragraph{Quién tiene el problema.} Personas que manejan su propio dinero —estudiantes universitarios con ingresos de mesadas o trabajos de medio tiempo, y trabajadores jóvenes con un salario fijo o ingresos variables— y que necesitan saber en qué se les va el dinero, si respetan los límites que se proponen y cuánto avanzan hacia una meta de ahorro.

\paragraph{Dónde ocurre.} En el manejo diario de las finanzas personales en Colombia. Nuestro contexto inmediato es la comunidad universitaria de Valledupar (Cesar), donde formulamos y probamos la solución.

\paragraph{Cómo ocurre hoy.} Quien intenta controlar sus finanzas lo hace de tres formas: lo apunta en un cuaderno, lleva una hoja de cálculo o confía en la memoria y en el extracto del banco. Las tres tienen el mismo problema: cada gasto se anota a mano, la suma por categoría y el porcentaje consumido de un presupuesto se calculan aparte, los gastos que se repiten (arriendo, suscripciones, salario quincenal) hay que volver a escribirlos cada vez y una meta de ahorro no se traduce en un plan concreto de cuánto guardar cada día. El extracto bancario solo muestra los movimientos de una cuenta y no los clasifica según los rubros de la persona. El resultado es que el registro se abandona o queda incompleto.

\paragraph{Evidencia.} La Encuesta de Medición de Capacidades Financieras de 2019, aplicada a 12\,000 hogares, encontró que solo el 51\,\% de los encuestados hace un presupuesto para manejar sus ingresos y gastos, y solo el 47\,\% anota sus gastos de manera regular \parencite{asobancaria2021}. Es decir, \textbf{más de la mitad de las personas (53\,\%) no lleva un registro regular de sus gastos} y casi la mitad (49\,\%) no tiene presupuesto. La misma fuente señala que el índice de educación financiera de Colombia retrocedió 1,13 puntos frente a la medición de 2013 de la CAF.

Para el contexto local, \pendiente{resultados de la encuesta o entrevista del equipo: número de personas, cómo registran hoy sus gastos, cuántas abandonaron el registro y por qué}.

\paragraph{Qué queda fuera del problema.} No abordamos la gestión de inversiones ni de créditos, la contabilidad de empresas o negocios, la declaración de impuestos, la conexión automática con cuentas bancarias ni el asesoramiento financiero. Tampoco pretendemos resolver la falta de educación financiera en sí: la reconocemos como causa, pero nuestra solución se dirige a la herramienta de registro y control (ver la sección \ref{sec:causas}).

\subsection{Causas y efectos}\label{sec:causas}

La Figura \ref{fig:arbol} resume el árbol de problemas. El problema central es que \emph{las personas no llevan un control sistemático de sus ingresos, gastos y metas de ahorro}.

\begin{figure}[H]
\centering
\caption{Árbol de problemas de ShiftCash}
\label{fig:arbol}
\resizebox{\linewidth}{!}{%
\begin{tikzpicture}[
  node distance=6mm and 4mm,
  efecto/.style={caja, text width=3.2cm, minimum height=1.3cm},
  causa/.style={caja, text width=3.2cm, minimum height=1.3cm},
  central/.style={caja, fill=acento, text=white, font=\small\bfseries, text width=13.6cm, minimum height=1cm},
]
\node[central] (P) {Las personas no llevan un control sistemático de sus ingresos, gastos y metas de ahorro};
% Efectos
\node[efecto, above=12mm of P.north west, anchor=south west] (E1) {E1. No saben en qué rubros se va su dinero};
\node[efecto, right=of E1] (E2) {E2. Superan los límites de gasto sin darse cuenta};
\node[efecto, right=of E2] (E3) {E3. Abandonan las metas de ahorro};
\node[efecto, right=of E3] (E4) {E4. Dificultad para cubrir obligaciones y ahorrar};
% Causas
\node[causa, below=12mm of P.south west, anchor=north west, cajaacento] (C1) {C1. Registro manual lento y propenso a errores u omisiones};
\node[causa, right=of C1, cajaacento] (C2) {C2. Sin cálculo automático del gasto por categoría ni del presupuesto};
\node[causa, right=of C2] (C3) {C3. Los gastos e ingresos recurrentes se olvidan};
\node[causa, right=of C3] (C4) {C4. Metas de ahorro sin un plan diario concreto};
\node[causa, below=6mm of $(C2.south)!0.5!(C3.south)$, text width=8cm, fill=grisclaro] (C5) {C5 (estructural). Bajo nivel de educación financiera};
\foreach \e in {E1,E2,E3,E4} \draw[flecha] (P.north -| \e.south) -- (\e.south);
\foreach \c in {C1,C2,C3,C4} \draw[flecha] (\c.north) -- (P.south -| \c.north);
\draw[flecha, gris] (C5.north) -- ++(0,3mm) -| (C1.south);
\draw[flecha, gris] (C5.north) -- ++(0,3mm) -| (C4.south);
\node[font=\footnotesize\itshape, anchor=east] at ([xshift=-2mm]E1.west) {Efectos};
\node[font=\footnotesize\itshape, anchor=east] at ([xshift=-2mm]C1.west) {Causas};
\end{tikzpicture}}
\figurenote{Las causas con borde de color son las que ataca ShiftCash como causa principal. Fuente: elaboración propia.}
\end{figure}

\begin{table}[H]
\caption{Causas del problema y hechos que las sustentan}
\label{tab:causas}
\begin{tabularx}{\linewidth}{L{0.8cm} L{4.2cm} Y}
\toprule
ID & Causa & Hecho que la sustenta \\
\midrule
C1 & El registro manual es lento y se omite & Solo el 47\,\% anota sus gastos con regularidad \parencite{asobancaria2021}; en cuaderno u hoja de cálculo cada movimiento se escribe a mano. \\
C2 & No hay cálculo automático del gasto ni del presupuesto & Saber cuánto se lleva gastado en una categoría exige sumar a mano; solo el 51\,\% hace presupuesto \parencite{asobancaria2021}. \\
C3 & Los movimientos recurrentes se olvidan & El arriendo, las suscripciones o el salario quincenal deben reescribirse en cada período. \\
C4 & Las metas no tienen plan & «Ahorrar un millón en tres meses» no dice cuánto guardar cada día ni cuánto se lleva. \\
C5 & Bajo nivel de educación financiera & El índice de educación financiera retrocedió 1,13 puntos frente a 2013 \parencite{asobancaria2021}. \\
\bottomrule
\end{tabularx}
\tablenote{Fuente: elaboración propia con datos de \textcite{asobancaria2021}.}
\end{table}

\begin{table}[H]
\caption{Efectos del problema}
\label{tab:efectos}
\begin{tabularx}{\linewidth}{L{0.8cm} L{4.2cm} Y}
\toprule
ID & Efecto & Relación con las causas \\
\midrule
E1 & No se sabe en qué rubros se va el dinero & Sin registro (C1) ni clasificación (C2) no hay totales por categoría. \\
E2 & Se superan los límites sin notarlo & Sin cálculo del porcentaje consumido (C2) el exceso se descubre a fin de mes. \\
E3 & Se abandonan las metas de ahorro & Sin plan diario ni seguimiento del avance (C4) la meta pierde sentido. \\
E4 & Dificultad para cubrir obligaciones y ahorrar & Es la consecuencia de E1 a E3 y de los gastos recurrentes olvidados (C3). \\
\bottomrule
\end{tabularx}
\tablenote{Fuente: elaboración propia.}
\end{table}

\paragraph{Causa que ataca ShiftCash.} La causa principal que atacamos es la \textbf{carga del registro manual sin cálculo automático} (C1 y C2): ShiftCash convierte cada movimiento en un registro de pocos segundos, lo clasifica por categoría y calcula solo los totales, el gasto de cada presupuesto y su porcentaje. Sobre esa base resuelve también C3 (las transacciones recurrentes se generan solas cada día) y C4 (cada meta recibe un plan de cuotas diarias y un indicador de progreso). La causa estructural C5 queda fuera del alcance: una aplicación no reemplaza un programa de educación financiera, aunque sí facilita el hábito de registrar.
